% !TEX root = ../glossario.tex
% Lettera R del Glossario (vedi capitoli/a.tex per le istruzioni).
\phantomsection
\letteraglossario{R}

\termine{RAG}{\emph{Retrieval-Augmented Generation}: tecnica in cui, prima di
  generare una risposta, si recuperano da una base di conoscenza i
  contenuti pertinenti alla domanda e li si fornisce a un \glo{LLM},
  così che la risposta si basi su fonti verificabili.}

\termine{Repository}{Archivio centralizzato in cui sono conservati i file del progetto, il codice sorgente e tutta la cronologia delle loro modifiche tramite un sistema di versionamento.}

\termine{Requisito}{Condizione o capacità che un prodotto software deve possedere per soddisfare un contratto, uno standard, una specifica o un altro documento formalmente imposto.}

\termine{Ruby on Rails}{Framework open source per lo sviluppo rapido di applicazioni web lato server (back-end), scritto nel linguaggio di programmazione Ruby e basato sul pattern architetturale Model-View-Controller.}